\begin{pastexam}[label={px:ust-109-ee-ds-05}]{臺聯大電機類 109 學年度 第 5 題}
（問答；10 分）

(10\%) Illustrate how data changes while using mergesort to sort: \{24, 4, 35, 1, 60, 12, 16, 52\}.
\end{pastexam}

\begin{pastsoln}
用資結書第~14.4~節由上而下的合併排序（虛擬碼~14.5）：從中間切成兩半，各自排好，再合併。
$8$ 個數切三層到只剩一個，再合併三層。

\begin{center}
\small
\begin{tabular}{ll}
\toprule
步驟 & 資料 \\
\midrule
原始 & $[24, 4, 35, 1, 60, 12, 16, 52]$ \\
分割 1 & $[24, 4, 35, 1]\quad[60, 12, 16, 52]$ \\
分割 2 & $[24, 4]\quad[35, 1]\quad[60, 12]\quad[16, 52]$ \\
分割 3 & $[24]\quad[4]\quad[35]\quad[1]\quad[60]\quad[12]\quad[16]\quad[52]$ \\
\midrule
合併 1 & $[4, 24]\quad[1, 35]\quad[12, 60]\quad[16, 52]$ \\
合併 2 & $[1, 4, 24, 35]\quad[12, 16, 52, 60]$ \\
合併 3 & $[1, 4, 12, 16, 24, 35, 52, 60]$ \\
\bottomrule
\end{tabular}
\end{center}

遞迴實際的執行順序是「先把左半做完再做右半」：$[24,4]\to[4,24]$，$[35,1]\to[1,35]$，合成
$[1,4,24,35]$；然後才處理右半 $[60,12]\to[12,60]$，$[16,52]\to[16,52]$，合成
$[12,16,52,60]$；最後合併兩半。 上表把同一層的結果放在一起，比較好看。

最後一次合併的過程（每次比較兩邊的第一個，小的先放；資結書虛擬碼~14.4）：
\begin{center}
\small
\begin{tabular}{lll}
\toprule
比較 & 放出 & 輸出 \\
\midrule
$1$ vs $12$ & $1$ & $1$ \\
$4$ vs $12$ & $4$ & $1, 4$ \\
$24$ vs $12$ & $12$ & $1, 4, 12$ \\
$24$ vs $16$ & $16$ & $1, 4, 12, 16$ \\
$24$ vs $52$ & $24$ & $1, 4, 12, 16, 24$ \\
$35$ vs $52$ & $35$ & $1, 4, 12, 16, 24, 35$ \\
左半用完 & $52, 60$ & $1, 4, 12, 16, 24, 35, 52, 60$ \\
\bottomrule
\end{tabular}
\end{center}
這次合併比較了 $6$ 次。 整個排序的比較次數：合併 1 每對 $1$ 次共 $4$ 次；合併 2 左邊
$[4,24]$ 和 $[1,35]$ 比 $3$ 次、右邊 $[12,60]$ 和 $[16,52]$ 比 $3$ 次；合併 3 比 $6$ 次；總共
$4 + 6 + 6 = 16$ 次。

$n = 8$ 是 $2$ 的次方，由下而上的合併排序（虛擬碼~14.6，一開始每個數自成一段，兩兩合併）切出來
的段和上表完全一樣，畫出來的過程相同。
\end{pastsoln}
